% ============================================================
% Section 3: Dataset & Methodology
% ============================================================

\section{Dataset \& Methodology}
\label{sec:methodology}

\subsection{AISHELL-5 Dataset}
\label{subsec:dataset}

AISHELL-5~\cite{aishell5_2025} is a Mandarin Chinese multi-speaker speech
corpus recorded in real vehicle cabins. Key characteristics:

\begin{itemize}
    \item \textbf{Recording setup}: Four directional microphones mounted at
    each car door (16\,kHz, 4-channel), yielding far-field recordings of
    natural cabin conversation.
    \item \textbf{Speakers}: 165 speakers, 2–4 per session, across 60+
    driving scenarios (highway, urban, parked).
    \item \textbf{Scale}: \textasciitilde100 hours total;
    $\sim$1.9\,GB \textit{dev} and $\sim$1.7\,GB \textit{eval1} subsets
    used for our experiments.
    \item \textbf{Supervision}: Precise per-speaker transcripts; near-field
    headset recordings (clean reference) available only in the training split.
    \item \textbf{License}: CC BY-SA 4.0 (OpenSLR \#159).
\end{itemize}

Dataset statistics for the \textit{dev} and \textit{eval1} splits used in
this work are summarized in Table~\ref{tab:dataset}.

\begin{table}[h]
\centering
\caption{AISHELL-5 subset statistics used in experiments.}
\label{tab:dataset}
\begin{tabular}{lccc}
\toprule
\textbf{Split} & \textbf{Size} & \textbf{\# Sessions} & \textbf{Avg. Duration}\\
\midrule
dev   & 1.9 GB & see EDA & see EDA \\
eval1 & 1.7 GB & see EDA & see EDA \\
eval2 & 1.8 GB & see EDA & see EDA \\
\bottomrule
\end{tabular}
\end{table}

\noindent\textit{Exact counts are populated by running \texttt{python evaluate.py --split dev --n -1}.}

\subsection{System Architecture}
\label{subsec:architecture}

Figure~\ref{fig:pipeline} illustrates the proposed five-stage pipeline.

\begin{figure}[h]
\centering
% \includegraphics[width=\columnwidth]{figures/system_diagram.png}
\fbox{\parbox{0.9\columnwidth}{\centering
\textbf{[Insert system\_diagram.png here]}\\
4-ch WAV $\to$ Separation $\to$ ASR $\to$\\
Speaker Role Classifier $\to$ Intent Engine}}
\caption{InCar-MSAR pipeline architecture. Four-channel audio
is chunked at 2\,s intervals, mixed to mono, separated into $N$ tracks
by SepFormer, transcribed by Whisper, role-classified by an energy-based
rule classifier (with optional ECAPA-TDNN), and parsed by a keyword
intent engine.}
\label{fig:pipeline}
\end{figure}

\subsubsection{Stage 1 – Multi-Channel Input}
Audio is loaded as a tensor of shape $[4, T]$ at 16\,kHz using
\texttt{torchaudio}. For streaming inference, the signal is segmented into
overlapping 2-second chunks (0.2\,s overlap-add) using a Hann window to
eliminate boundary artifacts.

\subsubsection{Stage 2 – Speech Separation}
Each chunk is mixed to mono ($[1, T]$) and processed by
SepFormer~\cite{subakan2021attention} (\texttt{speechbrain/sepformer-wsj02mix}),
yielding $N$ single-channel tracks. We configure $N\,{=}\,2$ for the majority
of sessions; the pipeline supports up to $N\,{=}\,4$. A fallback mechanism
bypasses separation and uses individual microphone channels directly when the
estimated separation quality (SI-SNRi proxy) falls below a 3\,dB threshold.

\subsubsection{Stage 3 – Automatic Speech Recognition}
Each separated track is transcribed using
Whisper-small~\cite{radford2023robust} (\texttt{openai/whisper-small},
$\approx$244\,M parameters) with \texttt{language="zh"}, beam size 2.
Silent chunks (RMS\,$<\,10^{-4}$) are detected and skipped to prevent
Whisper hallucination.

\subsubsection{Stage 4 – Speaker Role Classification}
An energy-based heuristic assigns the track most correlated with microphone
channel\,0 (driver-side door) the label \textit{Driver}; remaining tracks
are labeled \textit{Passenger\_N}. Optionally, ECAPA-TDNN~\cite{desplanques2020ecapa}
embeddings (\texttt{speechbrain/spkrec-ecapa-voxceleb}) are used for
session-level speaker clustering.

\subsubsection{Stage 5 – Intent Engine}
A keyword-matching engine maps Mandarin transcripts to 12 intent categories
(climate control, media, navigation, window, phone). Patterns are defined in
a human-readable YAML file and support entity extraction (e.g., temperature
values, music genres).

\subsection{Evaluation Metrics}
\label{subsec:metrics}

\begin{description}
    \item[WER] Per-speaker character error rate, computed as
    $\text{WER} = (S + D + I) / N_{\text{ref}}$ over Chinese characters
    (no word boundaries).
    \item[cpWER] Concatenated Permutation WER~\cite{watanabe2020chime6}:
    the minimum WER over all $N!$ permutations of hypothesis-to-reference
    speaker assignments, solved by the Hungarian algorithm.
    \item[Speaker Accuracy] Proportion of sessions where speaker roles
    are assigned correctly to all tracks.
    \item[Latency] Wall-clock time per 2-second chunk, measured on an
    NVIDIA T4 GPU (Docker, CUDA 11.8).
\end{description}
